\documentclass[../main.tex]{subfiles}

\begin{document}
% -----------------------------
\section{Scf.for}
% -----------------------------
In \texttt{MLIR} (and similarly in \texttt{xDSL}), \texttt{SCF} stands for Structured Control Flow, a dialect designed to represent high-level control flow
constructs such as loops and conditionals in a structured and analyzable form. Unlike low-level branch-based control flow,
\texttt{SCF} operations explicitly model structured constructs like if and for, making transformations and optimizations easier to reason about.

The scf.for operation represents a counted loop with a clear start value (lower bound), end value (upper bound),
and step. It also supports loop-carried variables, allowing values to be passed from one iteration to the next
in a well-defined \texttt{SSA} form. This makes scf.for particularly useful in ML compilers, where loops over tensor
dimensions and iteration spaces are common and need to remain analyzable during optimization and lowering.

We will implement a \texttt{pow} operation in our high-level dialect and then lower it into an scf.for loop
to express exponentiation using structured iteration.
Extend the HC dialect:
\begin{lstlisting} [language=Python, caption={Pow operation definition}]
@irdl_op_definition
class HCPow(IRDLOperation):
    name = "hc.pow"
    base = operand_def(IntegerType)
    exp  = operand_def(IntegerType)
    res  = result_def(IntegerType)

    def verify_(self):
        if self.base.type != self.exp.type or self.res.type != self.base.type:
            raise ValueError("hc.pow: base, exp, and res must have same integer type")
\end{lstlisting}
and register it:
\begin{lstlisting} [language=Python, caption={Registering HCPow}]
    (HCAdd, HCMul, HCSub, HCRelu, HCPow),      # register operations here
\end{lstlisting}
In \path{hc_lowering.py}, provide additional imports for \texttt{builtin}, \texttt{scf},
\texttt{Block} and \texttt{Region}. 
\begin{lstlisting} [language=Python, caption={Imports for lowering}]
from xdsl.dialects import arith, builtin, scf
from xdsl.ir import Block, Operation, Region
\end{lstlisting}
Extend the list of supported operations:
\begin{lstlisting} [language=Python, caption={Supported operations}]
if op.name not in ("hc.add", "hc.mul", "hc.sub", "hc.relu", "hc.pow"):
    return
\end{lstlisting}
and provide lowering logic:
\begin{lstlisting} [language=Python, caption={hc.pow lowering}]
# ---------------- hc.pow ----------------
if op.name == "hc.pow":
    base, exp = op.operands
    ty = op.results[0].type  # integer type

    # constants for loop bounds/step
    c0 = arith.ConstantOp.from_int_and_width(0, 32)
    c1 = arith.ConstantOp.from_int_and_width(1, 32)

    # init accumulator = 1
    init = arith.ConstantOp.from_int_and_width(1, 32)

    # scf.for (%i = 0; %i < exp; %i += 1) iter_args(%acc = init) -> (ty) { %acc2 = muli %acc, base; scf.yield %acc2 }
    body = Block(arg_types=[ty, ty])

    # Better: use index for iv and keep acc as ty.
    # Let's construct it robustly:
    iv_ty = builtin.IndexType()
    body = Block(arg_types=[iv_ty, ty])
    iv, acc = body.args

    mul = arith.MuliOp(acc, base)
    body.add_ops([mul, scf.YieldOp(mul.result)])

    loop = scf.ForOp(
        lb=c0.result,          # lower bound
        ub=exp,                # upper bound (dynamic)
        step=c1.result,        # step
        iter_args=[init.result],
        body=Region(body),
    )

    # scf.ForOp returns the iter_args results
    rewriter.replace_op(
        op,
        new_ops=[c0, c1, init, loop],
        new_results=[loop.results[0]],
        safe_erase=True,
    )
    return
\end{lstlisting}
Try it in our \path{hc_main.py}:
\begin{lstlisting} [language=Python, caption={}]
from hc_dialect import HiCompiler, HCAdd, HCSub, HCMul, HCRelu, HCPow
...
# Operand must be SSA value
c5 = arith.ConstantOp.from_int_and_width(3, 32)
hc_pow = HCPow(
    operands=[c2, c4], # Base and exponent
    result_types=[i32],
)
# Return the result
ret = func.ReturnOp(hc_pow.results[0])

# Block containing the constants + hc ops + return
block = Block(ops=[c1, c2, c3, c4, c5, hc_add, hc_mul, hc_sub, hc_relu, hc_pow, ret])
\end{lstlisting}

Lowered function should look something like (notice scf.for region-based operation):
\begin{lstlisting} [language=Python, caption={Lowered example}]
builtin.module {
  func.func @my_func() -> i32 {
    ...
    %10 = arith.constant 0 : i32
    %11 = arith.constant 1 : i32
    %12 = arith.constant 1 : i32
    %13 = scf.for %14 = %10 to %2 step %11 iter_args(%15 = %12) -> (i32) {
      %16 = arith.muli %15, %2 : i32
      scf.yield %16 : i32
    }
    func.return %13 : i32
  }
}
\end{lstlisting}
To support \texttt{pow} operation in the simulator, we need to extend \path{interpreter.py}:
\begin{lstlisting} [language=Python, caption={Extending interpreter with pow operation}]
from xdsl.ir import BlockArgument, Operation, SSAValue
...
            if name == "hc.pow":
                a, b = op.operands
                self._set(op.results[0], self._get(a) ** self._get(b))
                continue
...
            if name == "scf.for":
                # Try: recognize "pow lowering" pattern and shortcut it.
                if self._try_eval_pow_lowering(op):
                    continue
                raise RuntimeError("Unsupported scf.for (not recognized as pow lowering)")
\end{lstlisting}
And provide helper function used for \texttt{pow} pattern evaluation:
\begin{lstlisting} [language=Python, caption={Pow lowering recognition}]
    def _try_eval_pow_lowering(self, op) -> bool:
        """
        Recognize and evaluate the specific lowering:
          %res = scf.for %iv = %lb to %ub step %step iter_args(%acc = %init) -> (i32) {
            %m = arith.muli %acc, %base : i32
            scf.yield %m : i32
          }
        If matched: set scf.for result(s) and return True, else False.
        """
        # Must have exactly 1 iter_arg result
        if len(op.results) != 1:
            return False

        # Must have lb, ub, step at least
        if len(op.operands) < 4:
            return False

        lb = self._get(op.operands[0])
        ub = self._get(op.operands[1])
        step = self._get(op.operands[2])
        init = self._get(op.operands[3])  # first iter_arg init

        # Typical pow lowering: lb=0, step=1, ub>=0
        if lb != 0 or step != 1 or ub < 0:
            return False

        # Body: single block
        try:
            body_block = op.regions[0].blocks[0]
        except Exception:
            return False

        ops = list(body_block.ops)
        # Expect exactly: [arith.muli, scf.yield]
        if len(ops) != 2:
            return False

        mul, yld = ops
        if mul.name != "arith.muli" or yld.name != "scf.yield":
            return False

        # scf.yield must yield the mul result
        if len(yld.operands) != 1 or yld.operands[0] is not mul.results[0]:
            return False

        # mul operands should be: (iter_arg_block_arg, base_value)
        # We don't want to fully interpret block args; just identify the "base" operand:
        if len(mul.operands) != 2:
            return False

        a, b = mul.operands

        # In xDSL, block args are typically instances of BlockArgument.
        if isinstance(a, BlockArgument) and not isinstance(b, BlockArgument):
            base_val = self._get(b)
        elif isinstance(b, BlockArgument) and not isinstance(a, BlockArgument):
            base_val = self._get(a)
        else:
            # Either both are block args or both are not -> not the simple pow pattern
            return False

        # Evaluate
        result = init * (base_val ** ub)

        self._set(op.results[0], result)
        return True
\end{lstlisting}

\end{document}
